% ECONOMICS_PROBLEM: {"id": "F13-303", "title": "Commodity-trade fragmentation", "jel_code": "F13", "source_id": "OP-0348"}
% !TeX program = xelatex
\documentclass[11pt,a4paper]{article}
\usepackage[margin=23mm]{geometry}
\usepackage{fontspec}
\setmainfont{Latin Modern Roman}
\usepackage{amsmath,amssymb,mathtools}
\usepackage{xcolor,enumitem,booktabs,longtable,array,xurl,hyperref}
\definecolor{muted}{HTML}{53636C}
\hypersetup{colorlinks=true,urlcolor=blue,linkcolor=blue}
\setlength{\parindent}{0pt}
\setlength{\parskip}{5pt}
\newcommand{\R}{\mathbb R}
\newcommand{\E}{\mathbb E}
\newcommand{\N}{\mathbb N}
\newcommand{\ind}{\mathbf 1}
\newcommand{\Var}{\operatorname{Var}}
\newcommand{\Cov}{\operatorname{Cov}}
\newcommand{\supp}{\operatorname{supp}}
\DeclareMathOperator*{\argmin}{arg\,min}
\DeclareMathOperator*{\argmax}{arg\,max}
\newcommand{\entrymeta}[3]{{\sffamily\small\color{muted}\textbf{\texttt{#1}}\quad|\quad #2\par #3\par}}
\newcommand{\Problemref}[1]{\texttt{#1}}
\newcommand{\JELref}[2]{\texttt{#2} (JEL \texttt{#1})}
\begin{document}
\section*{Commodity-trade fragmentation}
\textbf{Problem ID:} \texttt{F13-303}\quad \textbf{JEL:} \texttt{F13}\par
% BEGIN ECONOMICS STATEMENT
\label{problem:OP-0348}
{\sffamily\small\color{muted}Original subject group: Trade, globalization and geoeconomics\par}
\label{definitions:problem:30007630}
\textbf{Audit classification:} Explicit published open question. The IMF conclusion explicitly proposes further work on food security, mineral transition, recycling and elasticities. The dynamic scarcity probability and general-equilibrium expansion are editorial.
\subsection*{Definitions and precise research target}
Fix countries and commodity markets for food, energy and critical minerals, with production capacities, inventories and substitution technologies. A trade-fragmentation intervention \(a\) restricts specified bilateral flows at stated dates. Prices, investment and consumption adjust in a declared general-equilibrium system. Let \(p_{kct}(a)\) be the real price of commodity \(k\) in country \(c\), \(q_{kct}(a)\) delivered quantity for the declared end use and \(W_c(a)\) household welfare. The task is to identify or bound dynamic changes in these outcomes and scarcity risk,
\[\Pr(q_{kc,t+h}(do(a))<\underline q_{kc}),\]
where \(\underline q_{kc}\) is an independently justified minimum requirement. Define the baseline policy path, commodity-grade aggregation and treatment of informal rerouting. Estimate substitution elasticities and supply-expansion lags rather than imposing instantaneous adaptation. Include links between energy prices, fertilizer production, food supply and mineral-intensive investment where empirically material. Distinguish temporary price redistribution from long-run productive-resource losses and evaluate impacts on countries with weak buffers. Report uncertainty from concentrated suppliers and correlated disruptions. Scenario results conditional on assumed blocs and elasticities do not establish the exact effects of an unspecified real-world fragmentation process.

\textbf{Additional formal definitions and scope (editorial).} Commodity index \(k\) fixes grade, physical units and geographic delivery definition; countries are indexed by \(c\), and \(h\in\mathbb N_0\). Inventories obey material-balance equations with storage/decay costs, and supply expansion has a lagged investment technology. Let \(q_{kc,t+h}(a)\) be delivered quantity for a declared end use: household food or energy consumption, or industrial input availability for minerals. Do not compare total industrial throughput with a household nutritional requirement. Scarcity risk is a probability under the intervention's structural law, with \(\underline q_{kc}>0\) fixed in the same physical, nutritional or engineering units as the chosen end-use quantity. Define \(W_c\) through household utility including expenditure and income effects; producer surplus alone does not identify national welfare. Partial-equilibrium independent commodity markets, linked multi-commodity production and fully general equilibrium are distinct assumption sets. Shock correlation and endogenous supplier switching must be stated before computing scarcity probabilities.

For the identification part, declare an admissible class \(\Theta\) of structural models, including preferences, technologies, shock laws, measurement/sampling rules and any equilibrium selector. If \(O\) denotes the complete observable history and \(T(\theta)\) the target defined above, the identified set is \(\mathcal I_T=\{T(\theta):\theta\in\Theta,\ \mathcal L_\theta(O)=\mathcal L_{\rm obs}(O)\}\); point identification means this set is a singleton. A \(do(a)\) comparison replaces stated policy/technology equations while fixing the declared exogenous law and re-solving the remaining equations. These definitions do not assert that an identification theorem holds without further restrictions.
\subsection*{Variants and assumption changes}
\begin{enumerate}
\item \textbf{Independent commodity markets (source subcase).} Reproduce a declared partial-equilibrium bloc scenario with commodity-specific demand/supply elasticities.
\item \textbf{Inventories, recycling and linked markets (source-motivated extension).} Add lagged capacity, secondary mineral supply and energy/fertilizer/food links; identify dynamic scarcity and household welfare effects.
\end{enumerate}
\subsection*{References and source audit}
\begin{enumerate}
\item Primary IMF WP/23/201 cover verifies six authors and October 2023 title/date; official working-paper DOI supplied. Locator: Section 10, printed pp.38--39 (PDF pp.39--40), final paragraph. Support: Explicit commodity extension agenda. Full PDF accessible. URL: \url{https://www.imf.org/-/media/files/publications/wp/2023/english/wpiea2023201-print-pdf.pdf}.
\end{enumerate}
\begin{enumerate}\raggedright
\item\hypertarget{definitions:source:30007630:1}{} Jorge A. Alvarez; Mehdi Benatiya Andaloussi; Chiara Maggi; Alexandre Sollaci; Martin Stuermer; Petia Topalova. 2023. \emph{Geoeconomic Fragmentation and Commodity Markets}. Section 10, printed pp.38--39 (PDF pp.39--40), final paragraph.
\url{https://www.imf.org/-/media/files/publications/wp/2023/english/wpiea2023201-print-pdf.pdf}.
DOI: \url{https://doi.org/10.5089/9798400252426.001}.
\end{enumerate}

\subsection*{Common conventions: Source mathematical conventions}

These conventions apply to each entry unless explicitly replaced.

\textbf{Models and identification.} A primitive model \(m\) specifies all probability laws, feasible choices, information, equilibrium and measurement rules used by its target. The admissible class \(\mathcal M\) is fixed independently of outcomes. If \(P_O(m)\) is its observable probability law and \(\tau(m)\) the target, its identified set at observable law \(P\) is
\[
 \mathcal I_\tau(P)=\{\tau(m):m\in\mathcal M,\ P_O(m)=P\}.
\]
Point identification means that this set is a singleton. An empty set signals incompatibility of data and assumptions. Coordinate bounds need not describe the joint set. Bounds are sharp only if every retained target is attainable or arbitrarily approachable by compatible models. Finite bounds require explicit restrictions on unobserved quantities; unrestricted confounding, hidden wealth, extrapolation or arbitrary equilibrium selection can yield uninformative sets.

\textbf{Causal interventions.} A potential outcome \(Y_i(p)\) is the outcome under a complete policy regime \(p\), including timing, financing, information and market spillovers. The target population and units are fixed before treatment. With interference, use the complete assignment vector or a justified exposure mapping; individual no-interference notation alone cannot identify an economy-wide intervention. A difference of mean potential outcomes does not require the cross-world joint distribution. Conditioning on post-treatment migration, survival, enrollment or employment changes the estimand and needs separate justification.

\textbf{Statistical statements.} An estimator, confidence set, forecast or test specifies a sampling law, model class and loss/coverage criterion. Uniform \((1-\alpha)\)-coverage means \(\inf_{m\in\mathcal M}\Pr_m\{\tau(m)\in C_n\}\geq1-\alpha\), or an explicitly asymptotic version. The whole parameter space trivially has coverage, so informativeness requires a stated width, risk or power objective. A forecasting improvement is relative to a named benchmark, information set and evaluation population.

\textbf{Equilibrium and optimization.} An equilibrium comprises feasible plans, optimality, beliefs where needed, budget constraints and all market/resource clearing. The primitives or a stated benchmark define each correspondence \(\mathcal E(p;m)\). Nonemptiness is assumed or proved, never inferred just from compact policy sets. A selected-equilibrium target uses a declared selection rule; a robust target quantifies over all permitted equilibria. Use suprema or infima unless attainment is established. An \(\varepsilon\)-optimal policy is feasible and within \(\varepsilon\) of the finite optimum value. Report an empty feasible set separately.

\textbf{Welfare and accounting.} Normative utility, interpersonal normalization, social weights, numeraire, discounting and the use of public receipts are primitives. Behavioral utility need not coincide with the welfare criterion. Household welfare includes ownership income and rents once; adding profits again to compensating variation generally double-counts them. A monetary equivalent is defined by an explicit transfer and price convention, with monotonicity and range conditions. Average effects, mechanism contributions, transfers and welfare losses are different targets. Infinite sums require convergence and dynamic feasible plans require terminal or no-Ponzi conditions.

\textbf{Source versions.} A working-paper date, revision date and journal publication year are distinguished. A DOI for a journal version is not automatically the DOI of an earlier manuscript. Locators refer to the stated version and printed pages where specified. A metadata-only check does not verify a claim about a full-text conclusion. Every entry's source subsection narrows the provenance claim accordingly.
% END ECONOMICS STATEMENT
\end{document}
